\documentclass[10pt,a4paper]{amsart}
\usepackage[margin=19mm]{geometry}
\usepackage{amsmath,amssymb,mathtools}
\usepackage[scr=rsfs,cal=euler]{mathalfa}
\usepackage{xfrac}
\usepackage[unicode,hidelinks,bookmarks=false]{hyperref}
\newcommand{\ie}{i.e.\ }
\newcommand{\cf}{cf.\ }
\newcommand{\resp}{respectively\ }
\newcommand{\Subsection}{\subsection}
\providecommand{\nobreakdash}{\nobreak\hbox{-}}
\setlength{\emergencystretch}{2em}
\multlinegap0pt
\usepackage{amssymb}
\usepackage{mathtools}
\newtheorem{lemma}{Lemma}
\newtheorem{proposition}{Proposition}
\newcommand{\N}{\mathbb N}
\title{Avalanches Thresholds in One-Dimensional Bak--Sneppen Models}
\author{Olivier H\'enard}
\date{Source: arXiv:2608.10992v1 (CC BY 4.0)}
\begin{document}
\maketitle

\noindent\emph{Source and licence.} Avalanches Thresholds in One-Dimensional Bak--Sneppen Models. Version arXiv:2608.10992v1 (CC BY 4.0), distributed by arXiv under the Creative Commons Attribution 4.0 International licence, http://creativecommons.org/licenses/by/4.0/, as recorded in the arXiv licence metadata for this exact version and shown on the abstract page https://arxiv.org/abs/2608.10992v1. Full licence notice: \texttt{licenses/arxiv-2608.10992-CC-BY-4.0.txt}.\par\smallskip

\noindent\textit{Editorial scope.} Counted unit: the article's proposition prop:isotropic-endpoint-domination (source0.tex lines 1564--1596). The article's complete original proof of it is delivered with it (source0.tex lines 1597--1733, one \texttt{proof} environment of 137 lines). No mathematics is written, repaired or re-derived: the statement and the proof are the article's own lines, in the article's own notation, and no retained line is substituted or re-flowed.  Retained uncounted as the source-local context the proof needs: lines 1042--1072; lines 1130--1152; lines 1434--1439.  Nothing was omitted from any retained span, so the package carries no omission record.  No retained line contains a citation, so the wrapper carries no bibliography and no bibliography entry.  No retained line contains a figure environment, an image-inclusion command or drawing code, so no image is delivered and the package declares no asset dependency.  Editorial formatting only, in addition: an AMS \texttt{amsart} preamble that restates the article's own declarations and macros used by the retained lines, namely the environments \texttt{lemma}, \texttt{proposition} on separate counters, together with the standard packages those lines need.  The retained environments print in this document as Lemma 1, Proposition 1 in order of appearance, whereas the article prints the same items with the numbering of its own full document.  The complete native source of the article as distributed in the arXiv e-print for this exact version is bundled unchanged as \texttt{sources/207376/source0.tex}.
\begin{lemma}[Tail flow from general initial data]
\label{lem:tail-flow-general-data}
Let $\mu$ be a probability law on $\N$. There exists a unique family
$q^\mu=(q_k^\mu)_{k\geq1}$ such that $q_1^\mu\equiv1$,
$q_k^\mu\in C^1([0,\infty))$ for $k\geq2$,
$q_k^\mu(0)=\mu([k,\infty))$, and
\begin{equation}
\label{eq:general-tail-flow}
        (q_k^\mu)'
        =
        \sum_{u=1}^{k-1}q_u^\mu
        (q_{k-u}^\mu-q_k^\mu)
        +q_k^\mu(1-q_k^\mu),
        \qquad k\geq2.
\end{equation}
For every $t\geq0$,
\[
        1=q_1^\mu(t)\geq q_2^\mu(t)\geq\cdots\geq0,
\]
and each $q_k^\mu$ is non-decreasing in $t$. Consequently,
$q^\mu(t)$ is the tail of a unique probability law $\mu_t$ on
$\N\cup\{\infty\}$, with
\[
        \mu_t(\{\infty\})
        =
        q_\infty^\mu(t)
        :=
        \lim_{k\to\infty}q_k^\mu(t),
\]
and $(\mu_t)_{t\geq0}$ is non-decreasing in stochastic order.
\end{lemma}

\begin{proposition}[Restarted one-sided sharpness]
\label{prop:restarted-one-sided}
Let $\mu$ be a probability law with finite support in $\N$, and let
$q^\mu$ be the tail flow of Lemma~\ref{lem:tail-flow-general-data}. Set
\[
        m_\mu(t):=\sum_{k\geq1}q_k^\mu(t),
        \qquad
        \tau_\mu:=\sup\{t\geq0:m_\mu(t)<\infty\}.
\]
Then
\begin{equation}
\label{eq:restarted-susceptibility-bound}
        \tau_\mu
        \leq
        2\log\left(1+\frac1{m_\mu(0)}\right),
\end{equation}
and
\begin{equation}
\label{eq:restarted-postcritical-explosion}
        q_\infty^\mu(\tau_\mu+s)>0,
        \qquad s>0.
\end{equation}
\end{proposition}

\begin{equation}
\label{eq:isotropic-endpoint-supersolution}
h_k(t)-h_k(s)
\geq
\int_s^t\Phi_k(h(u))\,du.
\end{equation}

\begin{proposition}[One-sided endpoint domination]
\label{prop:isotropic-endpoint-domination}
Fix $s\geq0$ with $\mathbb E[L_s]<\infty$ and $K\geq1$. Let
$q^{s,K}$ be the tail flow of Lemma~\ref{lem:tail-flow-general-data}
with initial data
\[
q_k^{s,K}(0)
:=
\mathbb P(L_s\wedge K\geq k).
\]
Then, for every $k\geq1$ and $t\geq0$,
\begin{equation}
\label{eq:isotropic-endpoint-domination}
q_k^{s,K}(t)
\leq
\mathbb P(L_{s+t}\geq k).
\end{equation}
Consequently, with $m_{s,K}:=\mathbb E[L_s\wedge K]$,
\begin{equation}
\label{eq:isotropic-explosion-upper}
t_{\infty}^{\mathrm{iso}}
\leq
s+2\log\left(1+\frac1{m_{s,K}}\right).
\end{equation}
In particular,
\begin{equation}
\label{eq:isotropic-one-sided-threshold-comparison}
X_t\preceq L_t,
\qquad
t_{\infty}^{\mathrm{iso}}
\leq t_{\infty}^{\mathrm{os}}.
\end{equation}
\end{proposition}

\begin{proof}
Set $\widetilde h_k(t):=h_k(s+t)$. Then
$(q_k^{s,K})'=\Phi_k(q^{s,K})$,
$q_k^{s,K}(0)\leq\widetilde h_k(0)$, and
\eqref{eq:isotropic-endpoint-supersolution} shows that
$\widetilde h_k$ is an integral supersolution of the same equation.
We prove $q_k^{s,K}\leq\widetilde h_k$ by induction on $k$. The assertion
is immediate for $k=1$.
Suppose it holds below level $k$ and set
$d_k:=q_k^{s,K}-\widetilde h_k$. For $0\leq r\leq t$,
\[
d_k(t)-d_k(r)
\leq
\int_r^t
\bigl(\Phi_k(q^{s,K}(u))-\Phi_k(\widetilde h(u))\bigr)\,du.
\]

For fixed \(x\in[0,1]\), consider
\[
        \Psi_x(a_1,\ldots,a_{k-1})
        :=
        \sum_{u=1}^{k-1}a_u(a_{k-u}-x)+x(1-x).
\]
For \(j<k\),
\[
        \partial_{a_j}\Psi_x
        =2a_{k-j}-x.
\]
Hence \(\Psi_x\) is non-decreasing in each coordinate on the region
where \(a_i\geq x\) for every \(i<k\).

By the induction hypothesis,
\(q^{s,K}_j(u)\leq\widetilde h_j(u)\) for \(j<k\), while
Lemma~\ref{lem:tail-flow-general-data} gives
\(q^{s,K}_j(u)\geq q^{s,K}_k(u)\). Therefore, with the \(k\)-th
coordinate fixed at \(q^{s,K}_k(u)\), increasing the lower coordinates
from \(q^{s,K}(u)\) to \(\widetilde h(u)\) can only increase
\(\Phi_k\). Consequently,
\[
\begin{aligned}
        \Phi_k(q^{s,K}(u))-\Phi_k(\widetilde h(u))
        &\leq
        \Phi_k\bigl(
        \widetilde h_1(u),\ldots,\widetilde h_{k-1}(u),
        q^{s,K}_k(u)
        \bigr)
        -\Phi_k(\widetilde h(u))  \\
        &=
        a_k(u)d_k(u),
\end{aligned}
\]
where
\[
        a_k(u)
        :=
        1-\sum_{j=1}^{k-1}\widetilde h_j(u)
        -q^{s,K}_k(u)-\widetilde h_k(u)
        \leq0,
\]
using $\tilde h_1 \equiv 1$. 
Thus
\begin{equation}
\label{eq:isotropic-endpoint-comparison-difference}
d_k(t)-d_k(r)
\leq
\int_r^t a_k(u)d_k(u)\,du.
\end{equation}
No continuity of $\widetilde h_k$ is needed here. Since
$\widetilde h_k$ is non-decreasing, every discontinuity of
$d_k=q_k^{s,K}-\widetilde h_k$ is downward; the integral inequality
\eqref{eq:isotropic-endpoint-comparison-difference} controls any possible
increase.
We claim that $d_k\leq0$. Otherwise, fix $t>0$ such that $d_k(t)>0$ and
let
\[
\tau:=\sup\{r\in[0,t]:d_k(r)\leq0\}.
\]
The set is nonempty since $d_k(0)\leq0$. On $(\tau,t]$ one has
$d_k>0$, and hence $a_kd_k\leq0$. Choose $r_n\uparrow\tau$ with
$d_k(r_n)\leq0$; if $\tau=0$, take $r_n=0$. Since $a_k$ and $d_k$ are
bounded on compact intervals, \eqref{eq:isotropic-endpoint-comparison-difference}
gives
\[
\begin{aligned}
d_k(t)
&\leq
d_k(r_n)
+\int_{r_n}^{\tau}|a_k(u)d_k(u)|\,du
+\int_{\tau}^{t}a_k(u)d_k(u)\,du \\
&\leq
\int_{r_n}^{\tau}|a_k(u)d_k(u)|\,du
\longrightarrow0.
\end{aligned}
\]
If $\tau=t$, the same estimate applies with the last integral absent.
In either case this contradicts $d_k(t)>0$. Thus $d_k\leq0$, completing
the induction and proving \eqref{eq:isotropic-endpoint-domination}.
Set
\[
\tau_{s,K}
:=
\sup\left\{t\geq0:
\sum_{k\geq1}q_k^{s,K}(t)<\infty\right\}.
\]
Proposition~\ref{prop:restarted-one-sided} gives
\[
\tau_{s,K}
\leq
2\log\left(1+\frac1{m_{s,K}}\right),
\qquad
q_\infty^{s,K}(\tau_{s,K}+\varepsilon)>0
\]
for every $\varepsilon>0$. Letting $k\to\infty$ in
\eqref{eq:isotropic-endpoint-domination} at time
$\tau_{s,K}+\varepsilon$ yields
\[
\mathbb P
(L_{s+\tau_{s,K}+\varepsilon}=\infty)
\geq
q_\infty^{s,K}(\tau_{s,K}+\varepsilon)>0.
\]
Hence
\[
t_{\infty}^{\mathrm{iso}}
\leq
s+2\log\left(1+\frac1{m_{s,K}}\right)+\varepsilon,
\]
and \eqref{eq:isotropic-explosion-upper} follows by letting
$\varepsilon\downarrow0$.
For $s=0$ and $K=1$, the initial data are
$q_1^{0,1}(0)=1$ and $q_k^{0,1}(0)=0$ for $k\geq2$. Uniqueness of the solution to the
triangular system therefore gives
$q_k^{0,1}(t)=\mathbb P(X_t\geq k)$. Thus
\eqref{eq:isotropic-endpoint-domination} gives $X_t\preceq L_t$, and
letting $k\to\infty$ gives
$t_{\infty}^{\mathrm{iso}}\leq t_{\infty}^{\mathrm{os}}$.
\end{proof}

\end{document}
